%!TEX root = ../main.tex
%----------------------------------------------------------------------
%  Appendix
%  Paper:   Theory-Informed Kernel Ridge Regression for Asset Pricing
%  Authors: Amedeo Andriollo and Juan F. Imbet
%----------------------------------------------------------------------

\section{Proofs}
\label{app:proofs}

%----------------------------------------------------------------------
\subsection{Proof of Proposition~\ref{prop:representer} (Representer Theorem)}
%----------------------------------------------------------------------

\begin{proof}
Decompose any $f \in \mathcal{H}$ as $f = f_{\parallel} + f_{\perp}$,
where $f_{\parallel}$ lies in the finite-dimensional subspace
$\mathcal{V} = \operatorname{span}\{k(\mathbf{X}_s, \cdot) : s = 1, \ldots, n\}$
and $f_{\perp} \in \mathcal{V}^{\perp}$.
By the reproducing property,
\[
f(\mathbf{X}_s)
= \langle f, k(\mathbf{X}_s, \cdot) \rangle_{\mathcal{H}}
= \langle f_{\parallel}, k(\mathbf{X}_s, \cdot) \rangle_{\mathcal{H}}
  + \langle f_{\perp}, k(\mathbf{X}_s, \cdot) \rangle_{\mathcal{H}}
= f_{\parallel}(\mathbf{X}_s),
\]
since $k(\mathbf{X}_s, \cdot) \in \mathcal{V}$ and
$f_{\perp} \perp \mathcal{V}$.
Therefore the prediction loss
$\frac{1}{n}\sum_s (R_s - f(\mathbf{X}_s))^2$
depends only on $f_{\parallel}$, and so does every $C_j(f)$
by assumption~\eqref{eq:penalty_eval}.

For the RKHS norm, orthogonality gives
$\|f\|_{\mathcal{H}}^2 = \|f_{\parallel}\|_{\mathcal{H}}^2 +
\|f_{\perp}\|_{\mathcal{H}}^2$.
Since $\|f_{\perp}\|_{\mathcal{H}}^2 \geq 0$ and
$\lambda_{\mathrm{stat}} > 0$, the objective is strictly decreased
by setting $f_{\perp} = 0$ whenever $f_{\perp} \neq 0$.
Hence any minimizer satisfies $f = f_{\parallel} \in \mathcal{V}$,
which yields the representation~\eqref{eq:representer}.
\end{proof}

%----------------------------------------------------------------------
\subsection{Proof of Proposition~\ref{prop:finite_dim} (Finite-Dimensional Problem)}
%----------------------------------------------------------------------

\begin{proof}
The prediction loss becomes
$\frac{1}{n}\|\mathbf{R} - \mathbf{K}\boldsymbol{\alpha}\|^2$
since $f(\mathbf{X}_s) = \sum_{s'} \alpha_{s'} k(\mathbf{X}_{s'}, \mathbf{X}_s)
= [\mathbf{K}\boldsymbol{\alpha}]_s$.
For the RKHS norm,
\[
\|f\|_{\mathcal{H}}^2
= \left\langle \sum_s \alpha_s\, k(\mathbf{X}_s, \cdot),\;
         \sum_{s'} \alpha_{s'}\, k(\mathbf{X}_{s'}, \cdot) \right\rangle_{\mathcal{H}}
= \sum_{s,s'} \alpha_s\, \alpha_{s'}\, k(\mathbf{X}_s, \mathbf{X}_{s'})
= \boldsymbol{\alpha}^\top \mathbf{K}\, \boldsymbol{\alpha}.
\]
By condition~\eqref{eq:penalty_eval}, each $C_j(f)$ is a function of
$\hat{\mathbf{f}} = \mathbf{K}\boldsymbol{\alpha}$.
Combining yields~\eqref{eq:finite_dim}.
\end{proof}

%----------------------------------------------------------------------
\subsection{Proof of Proposition~\ref{prop:foc} (First-Order Conditions)}
%----------------------------------------------------------------------

\begin{proof}
Write the objective~\eqref{eq:finite_dim} with quadratic penalties as
\begin{multline}
\label{eq:obj_expanded}
L(\boldsymbol{\alpha})
= \frac{1}{n}(\mathbf{R} - \mathbf{K}\boldsymbol{\alpha})^\top
  (\mathbf{R} - \mathbf{K}\boldsymbol{\alpha})
+ \lambda_{\mathrm{stat}}\,\boldsymbol{\alpha}^\top\mathbf{K}\boldsymbol{\alpha} \\
+ \sum_{j \in \mathcal{J}^A} \mu_j
  (\mathbf{A}_j\mathbf{K}\boldsymbol{\alpha} - \mathbf{b}_j)^\top
  (\mathbf{A}_j\mathbf{K}\boldsymbol{\alpha} - \mathbf{b}_j)
+ \sum_{j \in \mathcal{J}^{B}} \frac{\mu_j}{n}
  (\mathbf{K}\boldsymbol{\alpha} - \mathbf{v}_j)^\top
  (\mathbf{K}\boldsymbol{\alpha} - \mathbf{v}_j),
\end{multline}
where $\mathbf{v}_j = \mathbf{g}_j$ for Type~B.
Taking the gradient with respect to $\boldsymbol{\alpha}$:
\begin{align}
\frac{\partial L}{\partial \boldsymbol{\alpha}}
&= -\frac{2}{n}\,\mathbf{K}(\mathbf{R} - \mathbf{K}\boldsymbol{\alpha})
+ 2\lambda_{\mathrm{stat}}\,\mathbf{K}\boldsymbol{\alpha}
+ 2\sum_{j \in \mathcal{J}^A} \mu_j\,\mathbf{K}\mathbf{A}_j^\top
  (\mathbf{A}_j\mathbf{K}\boldsymbol{\alpha} - \mathbf{b}_j)
\notag \\
&\quad + 2\sum_{j \in \mathcal{J}^{B}} \frac{\mu_j}{n}\,\mathbf{K}
  (\mathbf{K}\boldsymbol{\alpha} - \mathbf{v}_j).
\label{eq:gradient}
\end{align}
Setting~\eqref{eq:gradient} to zero and left-multiplying by $\mathbf{K}^{-1}$
(which exists when $\mathbf{K} \succ 0$):
\[
\frac{1}{n}\,\mathbf{K}\boldsymbol{\alpha}
+ \lambda_{\mathrm{stat}}\,\boldsymbol{\alpha}
+ \sum_{j \in \mathcal{J}^A} \mu_j\,\mathbf{A}_j^\top\mathbf{A}_j
  \mathbf{K}\boldsymbol{\alpha}
+ \sum_{j \in \mathcal{J}^{B}} \frac{\mu_j}{n}\,\mathbf{K}\boldsymbol{\alpha}
= \frac{1}{n}\,\mathbf{R}
+ \sum_{j \in \mathcal{J}^A} \mu_j\,\mathbf{A}_j^\top\mathbf{b}_j
+ \sum_{j \in \mathcal{J}^{B}} \frac{\mu_j}{n}\,\mathbf{v}_j.
\]
Collecting terms using the definitions of $\boldsymbol{\Omega}$
and $\boldsymbol{\omega}$ from~\eqref{eq:Omega}--\eqref{eq:omega_vec}
gives~\eqref{eq:foc}.
Strict positive definiteness of $\mathbf{K}$ and $\lambda_{\mathrm{stat}} > 0$
ensure the coefficient matrix in~\eqref{eq:foc} is invertible, yielding
uniqueness.
\end{proof}

%----------------------------------------------------------------------
\subsection{Proof of Proposition~\ref{prop:nesting} (Nesting)}
%----------------------------------------------------------------------

\begin{proof}
Claims~(i) and~(ii) follow directly from the definition.
For~(i), setting $\mu_j = 0$ eliminates the structural penalties.
For~(ii), taking $\mu_j \to \infty$ enforces $C_j(f) = 0$ as a constraint;
formally, the sequence of penalized minimizers converges to the
constrained minimizer provided $C_j$ is lower semicontinuous and
$\{f \in \mathcal{H} : C_j(f) = 0\}$ is nonempty.
Setting $\lambda_{\mathrm{stat}} = 0$ removes agnostic regularization.
\end{proof}

%----------------------------------------------------------------------
\subsection{Proof of Proposition~\ref{prop:misspecification} (Misspecification Robustness)}
%----------------------------------------------------------------------

\begin{proof}
We show that any $\mu_j > 0$ introduces an asymptotic bias, so the
oracle sets $\mu_j^* = 0$.

\emph{Step 1: Asymptotic bias from misspecification.}
When $C_j(f^*) > 0$, the set
$\mathcal{F}_j(\mu_j) = \arg\min_{f \in \mathcal{H}}
\frac{1}{n}\sum_s (R_s - f(\mathbf{X}_s))^2
+ \lambda_{\mathrm{stat}}\|f\|_{\mathcal{H}}^2 + \mu_j C_j(f)$
converges (as $n \to \infty$ and $\lambda_{\mathrm{stat}} \to 0$ at an
appropriate rate) to the population minimizer
\[
f_j^*(\mu_j) = \arg\min_{f} \mathbb{E}[(R - f(\mathbf{X}))^2] + \mu_j C_j(f).
\]
By~\ref{cond:consistency}, $f_j^*(0) = f^*$.
For $\mu_j > 0$, the term $\mu_j C_j(f)$ tilts the solution toward
$\{f : C_j(f) \text{ is small}\}$.
Since $C_j(f^*) > 0$, this tilt is binding: $f_j^*(\mu_j) \neq f^*$
for all $\mu_j > 0$, and the prediction risk satisfies
\[
\mathbb{E}\bigl[(R - f_j^*(\mu_j)(\mathbf{X}))^2\bigr]
> \mathbb{E}\bigl[(R - f^*(\mathbf{X}))^2\bigr]
= \mathbb{E}[\varepsilon^2]
\qquad \text{for all } \mu_j > 0.
\]

\emph{Step 2: Risk gap is bounded away from zero.}
By continuity of $C_j$ (condition~\ref{cond:bounded}) and strict
convexity of the population prediction risk,
for any $\delta > 0$ there exists $c(\delta) > 0$ such that
\[
\mathrm{Risk}(\mu_j) - \mathrm{Risk}(0)
\;\geq\; c(\delta) > 0
\qquad \text{for all } \mu_j \geq \delta,
\]
where $\mathrm{Risk}(\mu_j) = \mathbb{E}[(R - f_j^*(\mu_j)(\mathbf{X}))^2]$.

\emph{Step 3: Oracle efficiency implies $\hat{\mu}_j \to 0$.}
By~\ref{cond:cv_oracle}, $\mathrm{CV}(\hat{\boldsymbol{\lambda}})$
tracks the oracle $\inf_{\boldsymbol{\lambda}} \mathrm{CV}(\boldsymbol{\lambda})$
up to $o_p(1)$.  Since the oracle selects $\mu_j^* = 0$
(the unique risk minimizer along the $\mu_j$ coordinate by Step~2),
and the risk gap is bounded below by $c(\delta)$ for $\mu_j \geq \delta$,
we have $\Pr(\hat{\mu}_j \geq \delta) \to 0$ for every $\delta > 0$.
\end{proof}

%----------------------------------------------------------------------
\subsection{Proof of Proposition~\ref{prop:test_mu} (Testing the Marginal Value of Economic Models)}
\label{app:proof_test_mu}
%----------------------------------------------------------------------

\begin{proof}
\emph{Under $H_0$.}
When $\mu_j^* = 0$, the oracle solution does not use model $j$.
Constraining $\mu_j = 0$ therefore imposes no binding restriction on the
population objective.
The difference
$\mathrm{CV}(\hat{\boldsymbol{\lambda}}_{-j}) -
\mathrm{CV}(\hat{\boldsymbol{\lambda}})$
arises solely from finite-sample noise in the cross-validated selection
of the remaining hyperparameters.
By condition~\ref{cond:test_rate} and the oracle efficiency
in~\ref{cond:cv_oracle}, this difference is
$O_p(n^{-1} r_n)$, so $T_j = n \cdot O_p(n^{-1} r_n) = O_p(r_n) = o_p(1)$.

\emph{Under $H_1$.}
When $\mu_j^* > 0$, constraining $\mu_j = 0$ removes a binding
restriction.
By condition~\ref{cond:test_smooth},
the population risk increases by at least
$c > 0$ (a positive constant depending on $\mu_j^*$ and $C_j$).
By uniform convergence of the cross-validated risk to the population risk,
$\mathrm{CV}(\hat{\boldsymbol{\lambda}}_{-j}) -
\mathrm{CV}(\hat{\boldsymbol{\lambda}}) \geq c/2 > 0$
with probability approaching one.
Hence $T_j = n(c/2 + o_p(1)) \to \infty$.
\end{proof}


%======================================================================
\section{Computational Implementation}
\label{app:computation}
%======================================================================

This appendix provides implementation details for reproducing the estimation results.
All computations are performed on a workstation with a 40-core CPU, 274~GB of RAM, and an NVIDIA RTX A6000 GPU with 48~GB of VRAM.

%----------------------------------------------------------------------
\subsection{Nystr\"{o}m Approximation}
\label{app:nystrom}
%----------------------------------------------------------------------

The full kernel matrix $\mathbf{K} \in \mathbb{R}^{n \times n}$ requires $O(n^2)$ storage and $O(n^3)$ time for the Cholesky factorization.
For the managed-portfolio specification ($n \approx 4{,}000$), the full kernel is feasible.
For the individual-stock panel ($n > 100{,}000$), the full kernel matrix exceeds 80~GB and is infeasible even on large-memory workstations.

We address this via the Nystr\"{o}m low-rank approximation.
Select $m \ll n$ landmark points uniformly at random from the training data and compute:
\begin{align}
\mathbf{K}_{nm} &\in \mathbb{R}^{n \times m}, \qquad [\mathbf{K}_{nm}]_{ij} = k(\mathbf{X}_i, \mathbf{X}_{l_j}), \label{eq:nystrom_knm} \\
\mathbf{K}_{mm} &\in \mathbb{R}^{m \times m}, \qquad [\mathbf{K}_{mm}]_{ij} = k(\mathbf{X}_{l_i}, \mathbf{X}_{l_j}), \label{eq:nystrom_kmm}
\end{align}
where $\{l_1, \ldots, l_m\}$ are the landmark indices.
The Nystr\"{o}m feature map is $\mathbf{Z} = \mathbf{K}_{nm}\,\mathbf{K}_{mm}^{-1/2} \in \mathbb{R}^{n \times m}$, where $\mathbf{K}_{mm}^{-1/2}$ is obtained via eigendecomposition.
The kernel matrix is then approximated as $\mathbf{K} \approx \mathbf{Z}\,\mathbf{Z}^\top$.

Under this approximation, the KRR solution becomes
\begin{equation}
\label{eq:nystrom_solve}
\hat{\boldsymbol{\beta}} = (\mathbf{Z}^\top\mathbf{Z} + n\lambda_{\mathrm{stat}}\,\mathbf{I}_m)^{-1}\,\mathbf{Z}^\top\mathbf{y},
\end{equation}
which requires only an $m \times m$ Cholesky factorization at cost $O(m^3)$, plus $O(nm)$ for the matrix-vector products.
Predictions on new data $\mathbf{X}_{\mathrm{new}}$ are computed as $\hat{\mathbf{f}} = \mathbf{Z}_{\mathrm{new}}\,\hat{\boldsymbol{\beta}}$, where $\mathbf{Z}_{\mathrm{new}} = \mathbf{K}_{\mathrm{new},m}\,\mathbf{K}_{mm}^{-1/2}$.

For Theory-KRR with structural penalties, the L-BFGS optimization operates on the $m$-dimensional coefficient vector $\boldsymbol{\beta}$ rather than the $n$-dimensional $\boldsymbol{\alpha}$.
Each L-BFGS iteration computes the fitted values $\hat{\mathbf{f}} = \mathbf{Z}\boldsymbol{\beta}$ at cost $O(nm)$, evaluates the structural penalties $C_j(\hat{\mathbf{f}})$, and updates $\boldsymbol{\beta}$.
The gradient with respect to $\boldsymbol{\beta}$ has a simple form:
\begin{equation}
\label{eq:nystrom_gradient}
\frac{\partial L}{\partial \boldsymbol{\beta}} = -2\,\mathbf{Z}^\top(\mathbf{y} - \mathbf{Z}\boldsymbol{\beta}) + 2n\lambda_{\mathrm{stat}}\,\boldsymbol{\beta} + \sum_{j} \mu_j\,\mathbf{Z}^\top\frac{\partial C_j}{\partial \hat{\mathbf{f}}},
\end{equation}
where $\partial C_j / \partial \hat{\mathbf{f}}$ is the $n$-dimensional gradient of the $j$-th penalty with respect to the fitted values.

We use $m = 500$ landmarks throughout.
Table~\ref{tab:nystrom_scaling} summarizes the computational scaling.

\begin{table}[h]
\centering
\caption{Computational Scaling: Full Kernel vs.\ Nystr\"{o}m Approximation}
\label{tab:nystrom_scaling}
\small
\begin{tabular}{lcc}
\toprule
 & Full Kernel & Nystr\"{o}m ($m = 500$) \\
\midrule
Kernel storage                & $O(n^2)$  & $O(nm)$ \\
KRR solve                     & $O(n^3)$  & $O(nm^2 + m^3)$ \\
L-BFGS per iteration          & $O(n^2)$  & $O(nm)$ \\
\midrule
$n = 4{,}000$ (portfolios)    & 128~MB, $<1$s & 16~MB, $<0.1$s \\
$n = 90{,}000$ (stocks)       & 65~GB, $\sim 10$h & 360~MB, $\sim 3$min \\
\bottomrule
\end{tabular}
\end{table}

%----------------------------------------------------------------------
\subsection{Coordinate Descent for Group Multipliers}
\label{app:coord_descent}
%----------------------------------------------------------------------

The hyperparameter vector $(\lambda_{\mathrm{stat}}, \mu_1^{\mathrm{group}}, \ldots, \mu_G^{\mathrm{group}})$ is $G + 1 = 9$ dimensional.
Exhaustive grid search with $k$ values per dimension requires $k^9$ evaluations; even with $k = 5$, this exceeds $10^6$ configurations, each requiring one full Theory-KRR fit.

We adopt cyclic coordinate descent over the group multipliers, following the algorithmic template of \citet{friedman2010regularization}.
The statistical penalty $\lambda_{\mathrm{stat}}$ is fixed by standard KRR cross-validation (a one-dimensional problem).
The remaining 8 group multipliers are optimized sequentially: for each group $g$, we evaluate 16 candidate values
\begin{equation}
\label{eq:cd_candidates}
\mu_g^{\mathrm{group}} \in \{0\} \cup \{10^{u} : u \in \mathrm{linspace}(-4, 1, 15)\},
\end{equation}
holding all other multipliers fixed, and retain the value that minimizes validation-set MSE.
Including zero in the candidate set allows the algorithm to shut off entire theory families when they do not improve prediction.

Each evaluation requires one Nystr\"{o}m KRR solve ($m \times m$ Cholesky, $m = 500$) plus a single Newton-like correction that incorporates the structural penalty gradient (see Section~\ref{app:nystrom}).
The 16 candidates within each group are evaluated in parallel across CPU cores, so the wall-clock cost per group equals a single evaluation.
One full cycle through $G = 8$ groups therefore costs $8 \times 1$ evaluations in wall-clock time; we run up to 3 cycles and stop early if no multiplier changes.

The total budget is at most $3 \times 8 \times 16 = 384$ evaluations, completing in under 5 seconds per rolling window on our 40-core workstation with Nystr\"{o}m features.
Coordinate descent has three advantages over random or grid search: (i)~it exploits the separable structure of the grouped penalty, (ii)~it automatically resolves substitution effects between correlated theory families, and (iii)~it scales linearly in the number of groups rather than exponentially.

%----------------------------------------------------------------------
\subsection{Parallelization Strategy}
\label{app:parallel}
%----------------------------------------------------------------------

We exploit three levels of parallelism:

\begin{enumerate}
\item \textbf{GPU acceleration.}
The kernel matrix computation~\eqref{eq:nystrom_knm}--\eqref{eq:nystrom_kmm} involves pairwise distance calculations and element-wise exponentials, which are trivially parallelizable on the GPU.
The pairwise distance computation for $n = 90{,}000$ points against $m = 500$ landmarks completes in under 2 seconds on the RTX A6000, compared to approximately 30 seconds on a single CPU core.
The subsequent $m \times m$ Cholesky factorization is performed on the CPU, as the $m = 500$ system is small enough that GPU transfer overhead dominates.

\item \textbf{Multi-threaded coordinate descent.}
Within each coordinate descent cycle, the 16 candidate values for a given group multiplier $\mu_g^{\mathrm{group}}$ are independent and evaluated in parallel across CPU threads.
With 40 threads, all 16 candidates complete in the wall-clock time of a single Nystr\"{o}m solve.
The sequential component---cycling through the $G = 8$ groups---is inherently serial, but each group's line search is fully parallel, yielding an effective speedup of $\min(16, 40) = 16\times$ relative to sequential evaluation.

\item \textbf{Parallel ablation studies.}
For predefined $\mu$ configurations (e.g., single-group activations used in the grid search baseline), each configuration is independent.
We distribute all configurations across $\min(C, 40)$ CPU threads, achieving near-linear speedup since they share no mutable state.
\end{enumerate}

The memory footprint scales linearly in the number of parallel workers, since each worker maintains its own copy of the Nystr\"{o}m feature matrix $\mathbf{Z}$.
With $m = 500$ and $n = 90{,}000$, each copy occupies approximately 360~MB, so 40 parallel workers require 14~GB---well within the 274~GB available on our workstation.
